\begin{textsample}{POS Dim 2 – ma – Score 27.00 – ma/ma\_em\_56.txt}  \label{ex:f2_pos_057}
5.4.
Sistemáticas de escolha, acompanhamento, avaliação e \textbf{mobilidade} no \textbf{itinerário} de formação \textbf{técnica} e \textbf{profissional} Os \textbf{itinerários} \textbf{formativos} correspondem a caminhos de formação compostos por unidades curriculares que visam ao aprofundamento e à ampliação de aprendizagens em uma ou mais áreas do conhecimento e/ou na formação \textbf{técnica} e \textbf{profissional} ( BRASIL, 2020 ).
Sua \textbf{implantação} na unidade escolar dar -se-á mediante escolha dos estudantes, que deverá ser orientada e apoiada pelos sistemas de ensino e pela comunidade escolar.
Nesse sentido, os sistemas de ensino precisam oferecer mais de um tipo de \textbf{itinerário} \textbf{formativo} em cada \textbf{município}, contemplando diversas áreas do conhecimento e/ou formação \textbf{técnica} e \textbf{profissional}, para proporcionar escolhas \textbf{formativas} aos estudantes de forma a \textbf{atender} à pluralidade de seus interesses e anseios, respeitando a heterogeneidade nas \textbf{ofertas} dos percursos \textbf{formativos} nas escolas de cada território.
Desse modo, a \textbf{oferta} dos aprofundamentos deve considerar, entre outros aspectos, os seguintes : - \textbf{Perfil} dos estudantes : os anseios e necessidades dos estudantes. - Número de estudantes : é recomendado que as escolas com maior número de estudantes ofereçam maior quantidade e variedade de aprofundamentos, objetivando \textbf{atender} à demanda de forma diversificada e ampla. - Equipe docente : a disponibilidade de tempo, as \textbf{habilitações} e \textbf{vocações} dos professores. - Infraestrutura : a estruturação e a quantidade dos espaços físicos, equipamentos e materiais existentes ou possíveis de serem adaptados ou adquiridos pela unidade de ensino, pelo sistema de ensino ou, ainda, por meio de parceria com outras instituições \textbf{ofertantes}. - Potencialidades locais : os potenciais, as demandas e particularidade do território em que a escola está localizada.
No caso específico da formação \textbf{técnica} e \textbf{profissional}, devem ser considerados, também, o potencial socioeconômico e ambiental, além das demandas do mercado de trabalho local e regional e das novas exigências ocupacionais ocasionadas pelas mudanças no mundo do trabalho.
O recomendável é que os estudantes tenham opções de escolhas quanto ao aprofundamento em, no mínimo, uma das quatro áreas do conhecimento ou formação \textbf{técnica} \textbf{profissional}, em atendimento às suas expectativas pessoais e \textbf{formativas}, embora em \textbf{itinerários} integrados, articulando mais de uma área do conhecimento e/ou formação \textbf{técnica} e \textbf{profissional}.
Os aprofundamentos podem ser \textbf{ofertados} na própria escola ou em unidades de ensino próximas e parceiras.
No que tange ao \textbf{itinerário} \textbf{formativo} \textbf{técnico} e \textbf{profissional}, foco desta proposta curricular, destaca -se : - A \textbf{carga} \textbf{horária} dos aprofundamentos varia conforme o \textbf{curso} \textbf{técnico}, o conjunto de FICs articuladas, incluindo ou não programa de aprendizagem \textbf{profissional} à estrutura curricular. - No caso de \textbf{cursos} \textbf{técnicos}, recomenda -se que, dentro da \textbf{carga} \textbf{horária}, seja considerado que o estudante possa percorrer formações articuladas à formação geral básica : diversificadas + básica para o trabalho + especificidades do \textbf{curso} + possibilidade de integração a outras áreas do conhecimento, conforme interesses dos estudantes e viabilidade de \textbf{oferta} pelo sistema de ensino.
Acompanhamento e avaliação Para atendimento dos estudantes da \textbf{rede} pública de ensino do Maranhão, considerará possibilidades de \textbf{ofertas} de \textbf{habilitações} e \textbf{qualificações} \textbf{técnicas} e \textbf{profissionais} e formação inicial continuada, presenciais ou a distância, contemplando as formas integrada, concomitante e subsequente, \textbf{primando} pelo \textbf{perfil} \textbf{profissional} de conclusão de cada \textbf{curso} \textbf{ofertado}, a partir da conexão com os eixos \textbf{estruturantes} e as possibilidades de \textbf{certificação} intermediária, continuada e verticalizada no \textbf{itinerário} \textbf{formativo}.
Nesse sentido, a participação de todos os envolvidos no processo educativo é substancial para as etapas de planejamento, acompanhamento e avaliação da \textbf{oferta}.
O que exigirá adoção de estratégias e instrumentos específicos de acompanhamento e avaliação dos processos de ensino e aprendizagem de modo sistemático e contínuo, objetivando a análise de dados educacionais para tomada de decisões coletivas quanto à necessidade de avanço, aprofundamento e reconhecimento de competências, experiências e saberes constituídos antes, durante e após o percurso \textbf{formativo}.
O processo avaliativo do \textbf{itinerário} \textbf{técnico} e \textbf{profissional} parte da concepção de que o estudante constitui o centro do processo ensino-aprendizagem.
O desenvolvimento integral, o protagonismo, o projeto de vida e a formação do estudante para conviver com os avanços e desafios do século XXI são basilares para estruturação e \textbf{efetivação} curricular ( BRASIL, 2020 ).
Para o monitoramento, o acompanhamento e a avaliação da \textbf{oferta} desse \textbf{itinerário} \textbf{formativo} executada tanto pelas próprias escolas quanto por instituições parceiras, faz -se necessária a criação de indicadores, diversificação de processos e instrumentos para avaliar a \textbf{qualidade} da \textbf{oferta} e o impacto dos percursos \textbf{formativos} na perspectiva do próprio estudante sobre a experiência pedagógica vivenciada, seu aprendizado e desenvolvimento das competências gerais da BNCC, das habilidades gerais relacionadas aos quatro eixos \textbf{estruturantes}, assim como das associadas ao mundo do trabalho, específicas de cada \textbf{habilitação} ou \textbf{qualificação} \textbf{técnica} e \textbf{profissional}.
A \textbf{certificação} Conforme Resolução nº 3/2018, para atendimento das exigências curriculares do ensino médio, os sistemas de ensino devem adotar critérios para reconhecimento de competências desenvolvidas pelos estudantes na formação geral básica e nos \textbf{itinerários} \textbf{formativos}, a partir das formas de comprovação, como avaliação de saberes ; demonstração prática ; e \textbf{documentação} expedida por instituições educativas.
No que tange ao \textbf{itinerário} \textbf{técnico} e \textbf{profissional}, as instituições e \textbf{redes} de ensino devem implementar processo de avaliação, reconhecimento e \textbf{certificação} de saberes e competências desenvolvidos no processo \textbf{formativo}, inclusive no ambiente de trabalho, com vistas a oportunizar \textbf{prosseguimento} ou conclusão de estudos, conferindo aos aptos um diploma, no caso de \textbf{curso} de \textbf{habilitação} \textbf{técnica} de nível médio, ou certificado, no caso de \textbf{curso} de \textbf{qualificação} \textbf{profissional} ( LDB nº 9.394/96 ). \textbf{Atendidas} as normas legais, as instituições e \textbf{redes} de ensino podem emitir certificado e/ou diploma ( \textbf{habilitação} \textbf{técnica} e \textbf{profissional} de nível médio ) de conclusão do ensino médio, validando os conhecimentos adquiridos pelos estudantes na formação geral básica e nos \textbf{itinerários} \textbf{formativos}.
Quando o aprofundamento do \textbf{itinerário} \textbf{técnico} e \textbf{profissional} ocorrer por instituições parcerias credenciadas e autorizadas para tal fim, a Resolução nº 3/2018 estabelece que a emissão de certificados e/ou diploma dar -se-á : - a instituição de ensino de origem do estudante é responsável pela expedição de certificado de conclusão do ensino médio ; - a instituição parceira deve expedir certificados, diplomas ou outros documentos comprobatórios das atividades pedagógicas concluídas sob sua responsabilidade ; - os certificados, diplomas ou outros documentos que comprovem atividades realizadas fora da escola de origem do estudante devem ser incorporados pela instituição de ensino de origem para efeito de \textbf{certificação} de conclusão do ensino médio ; - no caso da \textbf{habilitação} \textbf{técnica}, fica autorizada a instituição parceira a emitir diplomas de conclusão com validade apenas com a apresentação do certificado de conclusão do ensino médio.
Nos diplomas de \textbf{técnico} de nível médio, devem constar o título de \textbf{técnico} na respectiva \textbf{habilitação} \textbf{profissional}, o eixo tecnológico e/ou a área tecnológica cuja \textbf{habilitação} está vinculada, segundo o \textbf{Catálogo} Nacional de \textbf{Cursos} \textbf{Técnicos} ( CNCT ).
Já os certificados de \textbf{itinerários} compostos por um conjunto de FICs articuladas devem fazer referência à Classificação Brasileira de Ocupações ( CBO ).

% matched lemmas: atender, carga, catálogo, certificação, curso, documentação, efetivação, estruturante, formativo, habilitação, horário, implantação, itinerário, mobilidade, município, oferta, ofertante, ofertar, perfil, primar, profissional, prosseguimento, qualidade, qualificação, rede, técnico, vocação
\end{textsample}
